\section{Introduction}
\label{sec:introduction}

A heuristic specification inferrer for Java produces, for each method,
a set of Java Modeling Language (JML) clauses inferred from local
syntactic evidence: a null-dereference of a parameter inside the body
becomes a \texttt{requires \textbackslash p != null} clause; an array
read at index \texttt{i} becomes \texttt{requires 0 <= i \&\& i < a.length};
and so on. Recent work has shown that such specifications, even when
heuristically produced, materially improve downstream test
generation~\cite{edmonds2026inference} and can be embedded in compiled
artefacts with full lossless roundtrip~\cite{edmonds2026embedding}.
Both of those results depend on the spec set being \emph{rich enough}
to constrain downstream tools.

The dominant strategy for incorporating cross-method information in
such an inferrer is \emph{single-pass propagation}: while inferring
method \texttt{m}, the analyser looks up each callee's cached
specification, substitutes the actual arguments at the call site, and
appends the substituted callee preconditions to \texttt{m}'s
precondition set. This is fast and order-independent (provided the
callee is inferred first), and it is what the inferrer of
~\cite{edmonds2026inference} already does via its
\texttt{InterproceduralAnalyzer}. A natural question is whether this
single pass captures everything a \emph{genuinely compositional}
analysis would: that is, whether a second, top-down weakest-precondition
(WP) refinement pass, walking each method's body with the callee
specifications substituted at each call site, would discover any
additional clauses that the single pass missed.

The question matters operationally. If the answer is no, the
second-pass refactoring (and the associated implementation, testing,
and maintenance cost) is unjustified; the simpler single-pass
architecture is good enough. If the answer is yes, downstream
consumers of inferred specifications --- test generators, verifiers,
LLM-based tooling --- see a measurably richer spec set when the second
pass runs.

In this paper, we implement the second pass and answer the question
empirically on real Java code. We address the following research
questions:

\begin{itemize}\setlength\itemsep{0.2em}
  \item \textbf{RQ1.} Does a top-down WP-style compositional refinement
        pass strictly extend the precondition set produced by an
        established single-pass propagation, or is the single pass
        already sufficient?
  \item \textbf{RQ2.} If the second pass adds clauses, which structural
        shapes (null checks, range bounds, branch-conditional
        implications, disjunctive obligations) dominate the additions?
        The answer characterises \emph{why} the single pass falls
        short.
  \item \textbf{RQ3.} What is the runtime cost of the second pass at
        corpus scale, and how does it interact with cyclic
        strongly-connected components (mutual recursion) in the call
        graph?
\end{itemize}

We answer RQ1--RQ3 on Apache Commons Lang~3.14.0 (3{,}494 methods)
and Apache Commons IO~2.13.0 (2{,}111 methods), the same libraries
used by the predecessor work~\cite{edmonds2026inference,
edmonds2026embedding}. The principal findings are summarised in
Table~\ref{tab:headline} and analysed in Section~\ref{sec:empirical};
in short, the second pass refines 36--44\% of methods in the
specification cache and adds thousands of new precondition clauses,
with branch-conditional implications (\texttt{guard ==> obligation})
as the dominant shape.

\textbf{Contributions.} In summary, this paper makes the following
contributions:

\begin{itemize}\setlength\itemsep{0.2em}
  \item A 414-line \texttt{CompositionalAnalyzer} that implements the
        top-down WP refinement pass over a populated specification
        cache, supporting per-call-site argument substitution,
        enclosing-guard lifting, polymorphic-dispatch disjunction over
        candidate callees, and bounded fixpoint iteration on cyclic
        strongly-connected components.
  \item An empirical measurement on two real-world Java libraries
        (5{,}605 inferred methods in total) showing that the second
        pass refines 36--44\% of methods in the cache and adds 7{,}764
        new precondition clauses overall; the single-pass strategy is
        \emph{not} subsuming.
  \item A shape-categorisation of the newly-added preconditions
        showing that 53--57\% are branch-conditional implications
        (\texttt{guard ==> obligation}) --- the shape the single pass
        structurally cannot produce, because it lacks the control-flow
        context the WP walk supplies.
  \item Two real bugs found in the analyser by running it at scale: a
        \texttt{lastIndexOf('.')} in the polymorphic-candidate lookup
        that failed on parameter types containing dots
        (e.g.\ \texttt{java.lang.String}), and a related off-by-one
        in the call-site resolver. Both are fixed in the artefact
        accompanying this paper.
\end{itemize}

\textbf{Paper organisation.} Section~\ref{sec:background} reviews
single-pass propagation and the gap it leaves. Section~\ref{sec:algorithm}
describes the second-pass algorithm precisely.
Section~\ref{sec:impl} describes the implementation.
Section~\ref{sec:empirical} reports the measurement.
Section~\ref{sec:discussion} discusses what the result means for
single-pass inferrer architectures and for the downstream tools that
consume inferred specifications. Section~\ref{sec:threats} addresses
the threats to validity, Section~\ref{sec:related} positions the work
against related literature, and Section~\ref{sec:conclusion} concludes.
